\documentclass[10pt,a4paper]{extarticle}
\usepackage{parskip}

\usepackage[utf8]{inputenc}
\usepackage[T1]{fontenc}
\usepackage[brazil]{babel}
\usepackage{amsmath, amssymb, amsthm}
\usepackage{geometry}
\usepackage{listings}
\usepackage{xcolor}
\usepackage{hyperref}
\usepackage{booktabs}
\usepackage{array}
\usepackage{tabularx}
\usepackage{enumitem}
\usepackage{microtype}
\usepackage{csquotes}
\usepackage{tikz}
\usetikzlibrary{arrows.meta, positioning}

\definecolor{important}{RGB}{255, 100, 100}
\definecolor{notice}{RGB}{0, 102, 204}

\newcommand{\impt}[1]{\textcolor{important}{#1}}
\newcommand{\ntc}[1]{\textcolor{notice}{#1}}

\setlist{nosep}
\geometry{margin=2.2cm}

% Python listings
\definecolor{codebg}{rgb}{0.96,0.96,0.96}
\lstset{
    language=Python,
    backgroundcolor=\color{codebg},
    basicstyle=\ttfamily\footnotesize,
    keywordstyle=\color{blue},
    commentstyle=\color{gray},
    stringstyle=\color{red!70!black},
    showstringspaces=false,
    frame=single,
    breaklines=true,
    numbers=left,
    numberstyle=\tiny\color{gray},
    inputencoding=utf8
}

\newcommand{\pontos}[1]{\hfill\textnormal{\small\ntc{(#1 pontos)}}}

\begin{document}

\begin{center}
    {\LARGE\bfseries Avaliação (Versão B) --- Grafos, Algoritmos de Busca e NP-Completude}\\[0.3cm]
    {\large PCC104 --- Projeto e Análise de Algoritmos}\\[0.15cm]
    {DECOM --- ICEB --- UFOP}\\[0.35cm]
\end{center}

\noindent

\small
\textit{Instruções: Esta prova é individual e possui \textbf{5 questões}. É permitido consultar o código que você mesmo implementou nas listas de exercícios. Justifique todas as respostas --- respostas sem justificativa não recebem pontuação}

\normalsize
\vspace{0.4cm}

%============================================================
\begin{enumerate}[label=\textbf{\arabic*.}, itemsep=0.7cm, leftmargin=*]

%----------------------- Q1 -----------------------
\item \textbf{Busca em grafos.} \\
Considere a função de busca genérica cujo trecho principal é reproduzido abaixo. Cada vértice é armazenado como \texttt{[id, cor, pai, dist]}, com cores \texttt{'B'} (branco), \texttt{'C'} (cinza) e \texttt{'P'} (preto).

\begin{lstlisting}
while F:
    if type == 'dfs':
        x = F.pop(-1)
    elif type == 'bfs':
        x = F.pop(0)
    if x[0] in O:
        return x
    for xi in G[x[0]]:
        if V[xi][1] == 'B':
            V[xi][1] = 'C'
            V[xi][2] = x[0]
            V[xi][3] = V[x[0]][3] + 1
            F.append(V[xi])
    V[x[0]][1] = 'P'
\end{lstlisting}

\begin{enumerate}[label=(\alph*)]
    \item Descreva o significado das três cores no ciclo de vida de um vértice durante a busca. Em particular, em que instante do laço um vértice passa de \texttt{'C'} (cinza) para \texttt{'P'} (preto), e o que significa um vértice estar em \texttt{'C'}?
    \item O teste \texttt{if V[xi][1] == 'B'} controla quais vizinhos entram em \texttt{F}. Explique por que esse teste é essencial: o que aconteceria com a execução (e com a corretude) se ele fosse removido e todo vizinho fosse sempre inserido em \texttt{F}?
\end{enumerate}

%----------------------- Q2 -----------------------
\item \textbf{Representação de grafos.} \\
Considere um grafo com $n$ vértices representado por uma matriz de adjacência $M$ de dimensão $n \times n$.

\begin{enumerate}[label=(\alph*)]
    \item Descreva, em pseudocódigo ou em texto preciso, como implementar as operações \texttt{remover\_aresta(u,v)} e \texttt{inserir\_vertice(key)}. Indique a complexidade de cada uma em função de $n$.
    \item A mesma operação \texttt{inserir\_vertice(key)}, se o grafo estivesse representado por lista de adjacência, teria complexidade diferente? Compare os dois casos justificando.
\end{enumerate}

%----------------------- Q3 -----------------------
\item \textbf{Caminhos mínimos (Dijkstra).} \\
Considere o grafo direcionado ponderado abaixo, com origem no vértice $0$.

\begin{center}
\begin{tikzpicture}[
    >={Stealth[round]},
    node distance=2.6cm,
    vertex/.style={circle, draw, thick, minimum size=0.85cm, font=\bfseries},
    every edge/.style={draw, thick, ->}
]
\node[vertex] (0) {0};
\node[vertex] (1) [right=of 0] {1};
\node[vertex] (2) [below=2cm of 0] {2};
\node[vertex] (3) [below=2cm of 1] {3};
\node[vertex] (4) [right=of 3] {4};

\draw (0) edge node[above]{8} (1);
\draw (0) edge node[left]{4} (2);
\draw (1) edge node[right, pos=0.4]{2} (3);
\draw (2) edge node[pos=0.5, fill=white, inner sep=1pt]{2} (1);
\draw (2) edge node[above, pos=0.6]{7} (3);
\draw (2) edge[bend right=40] node[above]{3} (4);
\draw (3) edge node[above]{5} (4);
\draw (4) edge[bend left=30] node[above]{1} (3);
\end{tikzpicture}
\end{center}

\begin{enumerate}[label=(\alph*)]
    \item Execute o algoritmo de Dijkstra a partir de $0$. Apresente uma tabela mostrando, a cada extração de vértice, a distância \texttt{d} e o predecessor \texttt{pai} de cada vértice. Indique a distância final até todos os vértices.
    \item Na sua implementação (\texttt{dijkstra.py}), a atualização só ocorre quando \texttt{v['d'] > u['d'] + w[(u,v)]}. Explique o que essa condição representa e por que ela é essencial para a corretude do algoritmo.
\end{enumerate}

%----------------------- Q4 -----------------------
\item \textbf{NP-Completude.} \\
Um estudante quer provar que \textsc{Independent-Set} é NP-difícil e afirma: \textit{``Vou reduzir \textsc{Independent-Set} a \textsc{Clique}, que já sabemos ser NP-completo, e isso prova que \textsc{Independent-Set} é NP-difícil.''} O argumento está correto? Aponte o erro (se houver) e indique qual é a direção correta da redução entre \textsc{Clique} e \textsc{Independent-Set}.

%----------------------- Q5 -----------------------
\item \textbf{Vivendo com problemas NP-completos.} \\
Um laboratório de computação precisa resolver, rotineiramente, instâncias de um problema NP-completo de escalonamento de tarefas em máquinas. Diante da impossibilidade de uma garantia polinomial de pior caso, a equipe adota abordagens diferentes conforme o contexto. Para cada situação abaixo, identifique qual das cinco estratégias da Seção 8 das notas está sendo empregada (algoritmo exponencial rápido na prática, caso especial tratável, algoritmo de aproximação, heurística/metaheurística, ou complexidade parametrizada) e justifique em uma frase.

\begin{enumerate}[label=(\alph*)]
    \item Para uma estimativa rápida de tempo total, a equipe aceita uma alocação comprovadamente no máximo $50\%$ pior que a ótima, produzida por um procedimento garantidamente polinomial.
    \item Um cenário específico envolve tarefas cujas dependências formam sempre uma árvore; a equipe aplica um algoritmo próprio que resolve esse formato em tempo polinomial.
    \item Para o escalonamento noturno completo, roda-se um solver de \textit{branch-and-bound} que termina em segundos na esmagadora maioria dos casos, ainda que o pior caso permaneça exponencial.
    \item Um subproblema envolve apenas até $4$ máquinas críticas; usa-se um algoritmo cujo tempo é exponencial somente nesse número, mas linear no número total de tarefas.
    \item Sem qualquer garantia formal de qualidade, a equipe emprega um algoritmo genético que evolui as alocações e devolve a melhor solução encontrada dentro do orçamento de tempo.
\end{enumerate}

\end{enumerate}

\vfill
\begin{center}
\small\ntc{Boa prova!}
\end{center}

\end{document}
